\documentclass[11pt]{beamer}
\usetheme{Madrid}

\usepackage[T1]{fontenc}
\usepackage[utf8]{inputenc}
\usepackage{amsmath,amsfonts,amssymb}
\usepackage{graphicx}
\usepackage{xcolor}
\usepackage{tikz}
\usepackage{array}
\usepackage{longtable}

\title{RFE Response Toolkit}
\subtitle{EB2-NIW \& EB1A: Endeavor • 3 Prongs • GenAI • LaTeX • Pre-Filing Strategy}
\author{Strategic RFE Response System}
\date{\today}

\setbeamercovered{transparent}
\setbeamertemplate{navigation symbols}{}

\setbeamerfont{itemize/enumerate body}{size=\scriptsize}
\setbeamerfont{itemize/enumerate subbody}{size=\tiny}
\setbeamerfont{block body}{size=\scriptsize}
\setbeamerfont{block title}{size=\normalsize}

\begin{document}
	
	% TITLE SLIDE
	\begin{frame}
		\titlepage
		\centering
		\scriptsize \textit{Master the Art of RFE Response: EB2-NIW \& EB1A}
	\end{frame}
	
	% ============================================================
	% SECTION 1: INTRODUCTION
	% ============================================================
	
	\begin{frame}{Course Overview}
		\scriptsize
		\begin{block}{What You Will Learn}
			\begin{itemize}
				\item How to understand and respond to USCIS Requests for Evidence (RFEs)
				\item The 3 Prongs of EB2-NIW under \textit{Matter of Dhanasar}
				\item Using Generative AI (DeepSeek, ChatGPT, Gemini, Perplexity) for RFE responses
				\item LaTeX templates for professional 400+ page applications
				\item Expert opinion letters that address specific RFE concerns
				\item How to prove "Critical Role" for EB1A
				\item Aligning your research with U.S. National Interest
				\item \textcolor{blue}{NEW: Pre-filing work and pending publications strategy}
			\end{itemize}
		\end{block}
	\end{frame}
	
	\begin{frame}{What is an RFE?}
		\scriptsize
		\begin{block}{Request for Evidence - Definition}
			\textbf{Answer:} An RFE is a formal request from USCIS asking for additional documentation or clarification to determine eligibility.
			\vspace{0.2cm}
			\textbf{Key Facts:}
			\begin{itemize}
				\item Issued when evidence is missing, insufficient, or unclear
				\item Response time: typically 30-87 days (depending on the RFE)
				\item NOT a denial - your case is still alive
				\item Opportunity to strengthen your petition
			\end{itemize}
			\textcolor{red}{Do NOT ignore an RFE. Failure to respond = automatic denial.}
		\end{block}
	\end{frame}
	
	\begin{frame}{RFE vs. NOID vs. Denial}
		\scriptsize
		\begin{block}{Understanding USCIS Actions}
			\centering
			\begin{tabular}{|p{3cm}|p{5cm}|p{2.5cm}|}
				\hline
				\textbf{Action} & \textbf{Meaning} & \textbf{Response Window} \\
				\hline
				RFE & Core eligibility gaps found; case still active & 30-87 days \\
				\hline
				NOID & Intent to deny; serious defects & 30-33 days \\
				\hline
				Denial & Case closed; no further evidence accepted & 30-33 days for appeal \\
				\hline
			\end{tabular}
			\vspace{0.2cm}
			\textcolor{blue}{RFE is the best scenario - you have time to fix the issues.}
		\end{block}
	\end{frame}
	
	\begin{frame}{Why RFEs Happen - Common Triggers}
		\scriptsize
		\begin{block}{Most Common RFE Reasons for EB2-NIW/EB1A}
			\begin{enumerate}
				\item \textbf{Endeavor not clearly described} - officer cannot understand what you propose to do
				\item \textbf{National importance not established} - endeavor seems too local or employer-specific
				\item \textbf{Well-positioned not proven} - lack of evidence of progress or capability
				\item \textbf{Critical role not demonstrated} (EB1A) - role seems replaceable
				\item \textbf{Citation/impact evidence insufficient} - officer questions influence
				\item \textbf{Expert letters too generic} - not specifically addressing Dhanasar prongs
			\end{enumerate}
			\textcolor{blue}{Most RFEs are fixable with the right strategy.}
		\end{block}
	\end{frame}
	
	% ============================================================
	% SECTION 2: THE PRE-FILING STRATEGY (NEW ANGLE)
	% ============================================================
	
	\begin{frame}{The Pre-Filing Strategy: Work Done Before Filing}
		\scriptsize
		\begin{block}{Key Argument: Pending Publications and Work in Progress}
			\textbf{The Problem with RFEs:}
			\begin{itemize}
				\item USCIS often demands evidence of "already achieved" impact
				\item But the NIW standard is PROSPECTIVE (\textit{Dhanasar})
				\item Officers confuse EB-1A (past achievements) with EB2-NIW (future potential)
			\end{itemize}
			\textbf{The Strategic Answer:}
			\begin{itemize}
				\item \textcolor{blue}{Show that work was ALREADY IN PROGRESS before filing}
				\item \textcolor{blue}{Pending publications count as "work done" not "new achievements"}
				\item \textcolor{blue}{The papers were told/informed/described in the original petition}
				\item \textcolor{blue}{The research was ongoing - not started after filing}
			\end{itemize}
		\end{block}
	\end{frame}
	
	\begin{frame}{"The Papers Were Already Told" - Documentation Strategy}
		\scriptsize
		\begin{block}{Proving Pre-Filing Work Existed}
			\textbf{What you need to show:}
			\begin{itemize}
				\item The original petition mentioned the research direction
				\item Draft manuscripts existed BEFORE the filing date
				\item Preprints were submitted to arXiv/SSRN before filing
				\item Conference submissions were made before filing
				\item Research collaborations were established before filing
				\item Data collection was completed before filing
			\end{itemize}
			\textbf{How to prove it:}
			\begin{itemize}
				\item Email timestamps showing manuscript drafts
				\item Submission receipts from preprint servers (dated before filing)
				\item Conference acceptance letters (dated before filing)
				\item Version control history (GitHub, Overleaf) showing pre-filing work
				\item Expert letters confirming the work was underway
			\end{itemize}
			\textcolor{red}{If you can prove work was ongoing, it is NOT a post-filing achievement.}
		\end{block}
	\end{frame}
	
	\begin{frame}{Pending Publications Are NOT Post-Filing Achievements}
		\scriptsize
		\begin{block}{The Key Distinction Under Matter of Katigbak}
			\textbf{Katigbak bars:}
			\begin{itemize}
				\item New papers published AFTER the filing date
				\item New awards received AFTER the filing date
				\item New citations earned AFTER the filing date
			\end{itemize}
			\textbf{Katigbak does NOT bar:}
			\begin{itemize}
				\item \textcolor{green}{Papers that were SUBMITTED before filing (even if published later)}
				\item \textcolor{green}{Research that was IN PROGRESS at time of filing}
				\item \textcolor{green}{Draft manuscripts that existed before filing}
				\item \textcolor{green}{Work that was "already told" in the original petition}
			\end{itemize}
			\textbf{Strategic argument:}
			\begin{quote}
				``The pending publications represent the natural fruition of work that was already described, in progress, and ongoing at the time of filing. They do not create new eligibility; they merely confirm the viability and execution of the originally proposed endeavor.''
			\end{quote}
		\end{block}
	\end{frame}
	
	\begin{frame}{Incremental Achievements vs. New Achievements}
		\scriptsize
		\begin{block}{What Counts as "Incremental" vs. "New"}
			\centering
			\begin{tabular}{|p{4cm}|p{5cm}|}
				\hline
				\textbf{Incremental (Allowed)} & \textbf{New (Barred by Katigbak)} \\
				\hline
				Paper accepted that was submitted pre-filing & Brand new paper started post-filing \\
				\hline
				Citation count increase on pre-filing work & Citation from entirely new paper \\
				\hline
				Conference presentation of pre-filing research & New conference not mentioned before \\
				\hline
				Grant continuation of pre-filing project & New grant in different field \\
				\hline
				Expert letter confirming pre-filing capability & Expert letter from new unknown source \\
				\hline
			\end{tabular}
			\vspace{0.2cm}
			\textcolor{blue}{Argue that your achievements are INCREMENTAL (proving ongoing work) not NEW (creating new eligibility).}
		\end{block}
	\end{frame}
	
	\begin{frame}{How to Frame Pending Publications in Your RFE Response}
		\scriptsize
		\begin{block}{Sample Language for RFE Response}
			\begin{quote}
				``The RFE suggests that pending publications constitute impermissible post-filing achievements under \textit{Matter of Katigbak}. This is a misapplication of the precedent.
				
				The original petition explicitly described the Petitioner's ongoing research on [TOPIC]. Attached as Exhibit A are email communications showing that draft manuscripts were completed and submitted to journals on [DATE], which is BEFORE the priority date of this petition (March 31, 2025).
				
				The pending publications do not create new eligibility. Rather, they confirm that the proposed endeavor was concrete, viable, and already in progress at the time of filing. The work was already told. The publications are merely the natural output of work that was already described.
				
				USCIS has recognized that post-filing evidence may be considered where it corroborates facts existing at the time of filing. \textit{Matter of Katigbak} (14 I\&N Dec. 45). Here, the pending publications directly corroborate the Petitioner's claim that the research was ongoing.''
			\end{quote}
		\end{block}
	\end{frame}
	
	\begin{frame}{Building Your Pre-Filing Evidence Package}
		\scriptsize
		\begin{block}{What to Submit with Your RFE Response}
			\textbf{For each pending publication, submit:}
			\begin{itemize}
				\item [$\square$] Email from journal confirming submission date (BEFORE filing)
				\item [$\square$] Draft manuscript with version history showing pre-filing work
				\item [$\square$] Screenshot of preprint server submission (arXiv/SSRN) dated pre-filing
				\item [$\square$] Excerpt from original petition mentioning this research direction
				\item [$\square$] Expert letter confirming the work was discussed with them pre-filing
				\item [$\square$] Research log or lab notebook showing work dates
				\item [$\square$] Overleaf/GitHub commit history showing work before filing
			\end{itemize}
			\textcolor{blue}{The more documentation you have, the harder it is for USCIS to reject your argument.}
		\end{block}
	\end{frame}
	
	% ============================================================
	% SECTION 3: UNDERSTANDING THE 3 PRONGS
	% ============================================================
	
	\begin{frame}{Prong 1: Substantial Merit \& National Importance}
		\scriptsize
		\begin{block}{What USCIS Requires}
			\textbf{Substantial Merit:}
			\begin{itemize}
				\item Economic, scientific, cultural, or other benefits to the U.S.
				\item Can be shown through industry impact, publications, citations
			\end{itemize}
			\textbf{National Importance:}
			\begin{itemize}
				\item Broader implications beyond a single employer or location
				\item Potential for nationwide impact (prospective standard)
				\item Not required to have already achieved nationwide impact
			\end{itemize}
			\textcolor{red}{Common RFE error: Officer demands "nationwide adoption" - this is EB-1A standard.}
		\end{block}
	\end{frame}
	
	\begin{frame}{Prong 1 RFE Response Strategy}
		\scriptsize
		\begin{block}{How to Fix National Importance Deficiencies}
			\textbf{Evidence to submit:}
			\begin{itemize}
				\item .gov citations (Federal Reserve, Science.gov, ERIC)
				\item .edu citations (Harvard, MIT, Stanford citing your work)
				\item Expert letters explaining why your endeavor has national implications
				\item Documentation of broader adoption beyond your employer
				\item Letters from industry stakeholders describing your impact
				\item Evidence that your work is used across state lines
			\end{itemize}
			\textcolor{blue}{Show that your endeavor's benefits extend beyond your immediate location.}
		\end{block}
	\end{frame}
	
	\begin{frame}{Prong 2: Well-Positioned to Advance the Endeavor}
		\scriptsize
		\begin{block}{What USCIS Requires}
			\textbf{Evidence of being well-positioned:}
			\begin{itemize}
				\item Education and training relevant to the endeavor
				\item Experience and track record of success
				\item Progress made toward the endeavor (publications, prototypes)
				\item Plan for executing the endeavor (timeline, milestones)
				\item Access to resources, collaborators, or funding
			\end{itemize}
			\textcolor{red}{Common RFE error: Officer demands employer sponsorship or formal job offer.}
		\end{block}
	\end{frame}
	
	\begin{frame}{Prong 2 RFE Response Strategy}
		\scriptsize
		\begin{block}{How to Fix Well-Positioned Deficiencies}
			\textbf{Evidence to submit:}
			\begin{itemize}
				\item Detailed 5-year roadmap with annual milestones
				\item Documentation of prior successes in similar work
				\item Expert letters confirming your capability
				\item Evidence of ongoing projects (pre-filing)
				\item Letters from collaborators or potential partners
				\item Conference presentations, invited talks, workshops
				\item Peer review certificates (shows recognition by peers)
			\end{itemize}
			\textcolor{blue}{Self-petitioners are explicitly permitted under 8 C.F.R. § 204.5(k)(4)(ii).}
		\end{block}
	\end{frame}
	
	\begin{frame}{Prong 3: Balancing Test - Waiver Benefits the U.S.}
		\scriptsize
		\begin{block}{What USCIS Requires}
			\textbf{Factors favoring a waiver:}
			\begin{itemize}
				\item The endeavor is so urgent that waiting for labor certification would harm U.S. interests
				\item The nature of work makes standard job testing impossible
				\item The petitioner has unique qualifications that cannot be easily replicated
			\end{itemize}
			\textcolor{red}{Common RFE error: Officer applies PERM labor certification standard.}
		\end{block}
	\end{frame}
	
	\begin{frame}{Prong 3 RFE Response Strategy}
		\scriptsize
		\begin{block}{How to Fix Balancing Test Deficiencies}
			\textbf{Evidence to submit:}
			\begin{itemize}
				\item Letters explaining why your work cannot be tested through PERM
				\item Evidence of unique qualifications (rare combination of skills)
				\item Documentation of urgency (time-sensitive projects)
				\item Industry letters confirming shortage of experts in your niche
				\item Evidence that job description would not capture your specific endeavor
			\end{itemize}
			\textcolor{blue}{This prong is the most subjective - expert letters are particularly powerful here.}
		\end{block}
	\end{frame}
	
	% ============================================================
	% SECTION 4: THE ENDEAVOR
	% ============================================================
	
	\begin{frame}{What is an "Endeavor"?}
		\scriptsize
		\begin{block}{Defining Your Proposed Endeavor}
			\textbf{Answer:} The endeavor is the specific activity you propose to undertake in the United States.
			\vspace{0.2cm}
			\textbf{Key principles:}
			\begin{itemize}
				\item Must be clearly described (officer must understand what you will do)
				\item Must be specific, not vague or generic
				\item Must be something you are already doing or can immediately start
				\item Can be broader than a single job description
				\item Self-petitioners can propose independent research or entrepreneurship
			\end{itemize}
			\textcolor{red}{Vague endeavors = automatic RFE or denial.}
		\end{block}
	\end{frame}
	
	\begin{frame}{How to Describe Your Endeavor}
		\scriptsize
		\begin{block}{Structuring Your Endeavor Description}
			\textbf{Template for endeavor description:}
			\vspace{0.2cm}
			\begin{quote}
				``I propose to [specific activity] that will [specific outcome] by [specific methods]. This endeavor builds on my [X years] of experience in [field].''
			\end{quote}
			\textbf{Example:}
			\begin{quote}
				``I propose to continue my research on AI-driven financial risk modeling, focusing on developing open-source tools for community banks and credit unions to detect fraud and ensure regulatory compliance. This endeavor will enhance financial stability across the U.S. banking system.''
			\end{quote}
			\textcolor{blue}{Make it concrete, measurable, and connected to U.S. national interests.}
		\end{block}
	\end{frame}
	
	\begin{frame}{Relating Your Work to the Endeavor}
		\scriptsize
		\begin{block}{Connecting Past Work to Future Endeavor}
			\textbf{How to show continuity:}
			\begin{itemize}
				\item Show that your proposed endeavor is a natural extension of your existing work
				\item Cite specific pre-filing achievements that demonstrate capability
				\item Use expert letters to explain the relationship
				\item Create a timeline showing past $\rightarrow$ present $\rightarrow$ future
				\item Show that you have already made progress (publications, prototypes, collaborations)
			\end{itemize}
			\textcolor{red}{Disconnect between past work and proposed endeavor = Prong 2 deficiency.}
		\end{block}
	\end{frame}
	
	% ============================================================
	% SECTION 5: USING LATEX FOR RFE RESPONSES
	% ============================================================
	
	\begin{frame}{Why Use LaTeX for RFE Responses}
		\scriptsize
		\begin{block}{LaTeX Advantages for Immigration Filings}
			\textbf{Benefits:}
			\begin{itemize}
				\item Professional formatting (law-firm quality)
				\item Automatic table of contents and cross-references
				\item Consistent exhibit numbering
				\item Easy to update and version control
				\item Handles 400+ page documents seamlessly
				\item Free (Overleaf, TeX Live)
			\end{itemize}
			\textbf{What to templatize:}
			\begin{itemize}
				\item RFE response briefs
				\item Covering letters addressing each deficiency
				\item Exhibit indexes and organization
				\item Expert opinion letters
			\end{itemize}
		\end{block}
	\end{frame}
	
	\begin{frame}{LaTeX RFE Response Structure}
		\scriptsize
		\begin{block}{Organizing Your RFE Response in LaTeX}
			\begin{enumerate}
				\item \textbf{Main file (master\_rfe.tex):} \texttt{\textbackslash input} for each section
				\item \textbf{Cover letter (covering.tex):} Executive summary addressing each RFE point
				\item \textbf{Response to each deficiency:} Separate section per RFE issue
				\item \textbf{Exhibits:} Organized as Exhibit A, B, C...
				\item \textbf{Expert letters:} Individual files \texttt{\textbackslash input}
				\item \textbf{Evidence tables:} Publications, citations, peer reviews
			\end{enumerate}
			\textcolor{blue}{Use \texttt{\textbackslash includeonly} to test specific sections while compiling.}
		\end{block}
	\end{frame}
	
	% ============================================================
	% SECTION 6: GENERATIVE AI FOR RFE RESPONSES
	% ============================================================
	
	\begin{frame}{GenAI Tools for RFE Responses}
		\scriptsize
		\begin{block}{Using AI to Draft RFE Responses}
			\textbf{Recommended AI Tools:}
			\begin{itemize}
				\item \textbf{DeepSeek}: Best for legal reasoning and citation finding
				\item \textbf{ChatGPT}: Good for drafting and editing
				\item \textbf{Gemini}: Useful for summarization
				\item \textbf{Perplexity}: Great for finding supporting citations
				\item \textbf{Copilot}: Integration with LaTeX editors
			\end{itemize}
			\textcolor{blue}{Use multiple AI tools - each has different strengths.}
		\end{block}
	\end{frame}
	
	\begin{frame}{AI Prompt Template for RFE Response}
		\scriptsize
		\begin{block}{Sample Prompt for Prong 1 RFE}
			\begin{quote}
				USCIS denied Prong 1 (national importance) of my EB2-NIW petition stating: [COPY RFE TEXT].
				
				My proposed endeavor is: [DESCRIBE ENDEAVOR].
				
				My evidence includes: [LIST EVIDENCE].
				
				Draft a response arguing that:
				
				1. The endeavor has broader implications beyond my employer
				2. The officer applied an impermissibly heightened standard
				3. My evidence demonstrates potential for nationwide impact
				
				Cite Matter of Dhanasar and the USCIS Policy Manual.
			\end{quote}
		\end{block}
	\end{frame}
	
	\begin{frame}{AI for Expert Opinion Letters}
		\scriptsize
		\begin{block}{Using AI to Draft Expert Letters for RFEs}
			\textbf{AI prompt for expert letter:}
			\begin{quote}
				Draft an expert opinion letter for Dr. [NAME], PhD in [FIELD], supporting an EB2-NIW petition.
				
				The proposed endeavor is: [DESCRIBE].
				
				The RFE challenges Prong 1 (national importance).
				
				The letter should:
				- State the expert's qualifications
				- Describe how they know the petitioner's work
				- Explain why the endeavor has national importance
				- Address the specific RFE concerns
				- Conclude that the petitioner meets Dhanasar standards
			\end{quote}
			\textcolor{red}{Always have experts review and sign AI-drafted letters.}
		\end{block}
	\end{frame}
	
	% ============================================================
	% SECTION 7: EXPERT OPINION LETTERS
	% ============================================================
	
	\begin{frame}{Why Expert Letters Are Critical for RFEs}
		\scriptsize
		\begin{block}{Expert Letters Address Specific Deficiencies}
			\textbf{What expert letters can fix:}
			\begin{itemize}
				\item Prong 1: Confirm national importance from independent perspective
				\item Prong 2: Validate petitioner's capability and unique qualifications
				\item Prong 3: Explain why standard labor certification is inappropriate
				\item Help overcome officer's skepticism
				\item Provide third-party validation of your claims
			\end{itemize}
			\textcolor{red}{Generic letters are ignored. Letters must address specific RFE concerns.}
		\end{block}
	\end{frame}
	
	\begin{frame}{Expert Letter Structure for RFE Response}
		\scriptsize
		\begin{block}{Template for RFE-Focused Expert Letter}
			\begin{enumerate}
				\item \textbf{Expert qualifications:} Credentials, experience, publications
				\item \textbf{Relationship to petitioner:} How they know your work
				\item \textbf{Understanding of the endeavor:} Restate your proposed work
				\item \textbf{Response to specific RFE issue \#1:} Directly address each deficiency
				\item \textbf{Response to specific RFE issue \#2:} Directly address each deficiency
				\item \textbf{Conclusion:} Statement that you meet Dhanasar standards
				\item \textbf{Signature block:} Dated, signed, contact information
			\end{enumerate}
			\textcolor{blue}{Each RFE deficiency should have its own section in the expert letter.}
		\end{block}
	\end{frame}
	
	\begin{frame}{8 Variations of Expert Opinion Letters}
		\scriptsize
		\begin{block}{Different Letter Types for Different RFEs}
			\textbf{Variations included in toolkit:}
			\begin{enumerate}
				\item \textbf{General NIW support letter} - broad endorsement
				\item \textbf{Prong 1 focused} - national importance only
				\item \textbf{Prong 2 focused} - well-positioned only
				\item \textbf{Prong 3 focused} - balancing test only
				\item \textbf{RFE response specific} - addresses officer's exact language
				\item \textbf{Projection/adoption focused} - future impact
				\item \textbf{Critical role (EB1A)} - demonstrates indispensability
				\item \textbf{Research profile} - validates scholarly impact
			\end{enumerate}
			\textcolor{blue}{Choose the variation that matches your specific RFE deficiency.}
		\end{block}
	\end{frame}
	
	% ============================================================
	% SECTION 8: CRITICAL ROLE FOR EB1A
	% ============================================================
	
	\begin{frame}{EB1A: Proving "Critical Role"}
		\scriptsize
		\begin{block}{What USCIS Requires for Critical Role}
			\textbf{Evidence of critical role:}
			\begin{itemize}
				\item You played a leading or essential role in distinguished organizations
				\item Your contributions were not easily replaceable
				\item The organization's success depended on your involvement
				\item Your role was more than just "normal job duties"
			\end{itemize}
			\textcolor{red}{RFE common error: Officer says your role was not "critical" or was "normal duties."}
		\end{block}
	\end{frame}
	
	\begin{frame}{Critical Role RFE Response Strategy}
		\scriptsize
		\begin{block}{How to Prove Critical Role}
			\textbf{Evidence to submit:}
			\begin{itemize}
				\item Letters from supervisors describing your unique contributions
				\item Documentation of projects that succeeded because of your work
				\item Evidence that you were the lead on key initiatives
				\item Comparison of your role vs. others at same level
				\item Metrics showing your impact (revenue saved, efficiency gained)
				\item Awards or recognition within the organization
				\item Evidence that replacing you would be difficult or costly
			\end{itemize}
			\textcolor{blue}{Focus on outcomes that would not have happened without you.}
		\end{block}
	\end{frame}
	
	% ============================================================
	% SECTION 9: RESEARCH PROFILE
	% ============================================================
	
	\begin{frame}{Aligning Research with U.S. National Interest}
		\scriptsize
		\begin{block}{Why Research Must Connect to National Priorities}
			\textbf{U.S. national interest areas (highly valued):}
			\begin{itemize}
				\item Artificial Intelligence and Machine Learning
				\item Cybersecurity and Critical Infrastructure
				\item Financial Stability and Risk Management
				\item Healthcare and Public Health
				\item Workforce Development and Education
				\item Energy and Climate Technology
			\end{itemize}
			\textcolor{red}{Research in these areas is more likely to satisfy Prong 1 national importance.}
		\end{block}
	\end{frame}
	
	\begin{frame}{Validating Journal Rankings}
		\scriptsize
		\begin{block}{Scopus, Web of Science, and EBSCO}
			\textbf{How to prove journal quality:}
			\begin{itemize}
				\item Check if journal is in Scopus (scopus.com)
				\item Check if journal is in Web of Science (clarivate.com)
				\item Check if journal is in EBSCO (used by U.S. Commerce Library)
				\item Avoid Beall's List of predatory journals
				\item Document journal acceptance rate and impact factor
			\end{itemize}
			\textbf{U.S. Commerce Library:}
			\begin{itemize}
				\item Indexes journals through EBSCO Business Source Complete
				\item Having your work in Commerce Library = .gov recognition
				\item Verifiable by USCIS through search.library.doc.gov
			\end{itemize}
		\end{block}
	\end{frame}
	
	% ============================================================
	% SECTION 10: NOID RESPONSE
	% ============================================================
	
	\begin{frame}{NOID: Notice of Intent to Deny}
		\scriptsize
		\begin{block}{What is a NOID?}
			\textbf{Answer:} A NOID is more serious than an RFE - USCIS intends to deny but gives you one last chance.
			\vspace{0.2cm}
			\textbf{Key differences from RFE:}
			\begin{itemize}
				\item Shorter response window (30-33 days)
				\item Officer has already decided to deny unless you prove otherwise
				\item Requires stronger evidence and legal arguments
				\item Often involves legal errors, not just missing evidence
			\end{itemize}
			\textcolor{red}{Treat a NOID as a near-denial. Respond aggressively with legal arguments.}
		\end{block}
	\end{frame}
	
	\begin{frame}{NOID Response Strategy}
		\scriptsize
		\begin{block}{How to Respond to a NOID}
			\textbf{Step-by-step:}
			\begin{enumerate}
				\item Identify the legal errors in the NOID (wrong standard, misapplied precedent)
				\item Gather evidence directly addressing each error
				\item Draft legal brief citing \textit{Dhanasar}, Policy Manual, and case law
				\item Include expert letters specifically rebutting the NOID's reasoning
				\item Use same LaTeX template as RFE but with stronger legal arguments
				\item Respond within 30 days (no extensions)
			\end{enumerate}
			\textcolor{blue}{A well-crafted NOID response can lead to approval.}
		\end{block}
	\end{frame}
	
	% ============================================================
	% SECTION 11: I-290B
	% ============================================================
	
	\begin{frame}{I-290B: Notice of Appeal or Motion}
		\scriptsize
		\begin{block}{What is Form I-290B?}
			\textbf{Answer:} Form used to appeal a denial or file a motion to reopen/reconsider.
			\vspace{0.2cm}
			\textbf{Options on I-290B:}
			\begin{itemize}
				\item \textbf{Appeal to AAO:} Review by Administrative Appeals Office (no new evidence)
				\item \textbf{Motion to Reopen (MTR):} New evidence not available at filing
				\item \textbf{Motion to Reconsider (MTC):} Legal error in original decision
				\item \textbf{Combined MTR/MTC:} Both new evidence and legal errors
			\end{itemize}
			\textcolor{red}{Deadline: 30 days from denial date. No extensions.}
		\end{block}
	\end{frame}
	
	% ============================================================
	% SECTION 12: CONCLUSION & CHECKLISTS
	% ============================================================
	
	\begin{frame}{Final RFE Response Checklist}
		\scriptsize
		\begin{block}{Before Submitting Your RFE Response}
			\begin{itemize}
				\item [$\square$] Each RFE deficiency addressed individually
				\item [$\square$] Legal arguments cite \textit{Dhanasar} and Policy Manual
				\item [$\square$] Expert letters are specific (not generic)
				\item [$\square$] Exhibits organized with clear indexing
				\item [$\square$] Response submitted within deadline
				\item [$\square$] Proof of service attached (if required)
				\item [$\square$] Copy retained for your records
				\item [$\square$] Delivery tracking obtained
				\item [$\square$] \textcolor{blue}{Pre-filing work documented with timestamps}
				\item [$\square$] \textcolor{blue}{Pending publications linked to original endeavor description}
			\end{itemize}
			\textcolor{red}{Missing any RFE point = automatic denial.}
		\end{block}
	\end{frame}
	
	\begin{frame}{Common RFE Response Mistakes to Avoid}
		\scriptsize
		\begin{block}{What NOT to Do}
			\begin{enumerate}
				\item \textbf{Responding late:} Even one day late = automatic denial
				\item \textbf{Missing an RFE point:} Address every single issue raised
				\item \textbf{Submitting same evidence:} Must add NEW evidence or arguments
				\item \textbf{Generic expert letters:} Must address specific RFE language
				\item \textbf{Poor organization:} Officer must find your response easily
				\item \textbf{Emotional language:} Stick to facts and law
				\item \textbf{Forgetting exhibits:} Evidence must support every claim
				\item \textbf{\textcolor{red}{Claiming post-filing achievements as new evidence}} (Katigbak violation)
			\end{enumerate}
			\textcolor{blue}{Organization is critical - make the officer's job easy.}
		\end{block}
	\end{frame}
	
	\begin{frame}{RFE Response Timeline}
		\scriptsize
		\begin{block}{Strategic Timeline for RFE Response}
			\begin{tabular}{|p{3cm}|p{6cm}|}
				\hline
				\textbf{Week} & \textbf{Action} \\
				\hline
				Week 1 & Read RFE carefully; identify each deficiency \\
				\hline
				Week 2 & Gather new evidence; contact experts for letters \\
				\hline
				Week 3 & Draft response using LaTeX template \\
				\hline
				Week 4 & AI-assisted refinement; expert review \\
				\hline
				Week 5 & Paralegal review; final formatting \\
				\hline
				Week 6 & Submit response (early, not on deadline day) \\
				\hline
			\end{tabular}
			\vspace{0.2cm}
			\textcolor{red}{Do NOT wait until the last day. Submit at least 1 week early.}
		\end{block}
	\end{frame}
	
	\begin{frame}{Resources and Templates Included}
		\scriptsize
		\begin{block}{What You Get in This Course}
			\textbf{LaTeX Templates:}
			\begin{itemize}
				\item Master RFE response template (50+ pages)
				\item Cover letter template
				\item Expert opinion letter (8 variations)
				\item Exhibit organization template
			\end{itemize}
			\textbf{AI Prompts:}
			\begin{itemize}
				\item Prompts for DeepSeek, ChatGPT, Gemini, Perplexity
				\item Prompts for each prong and each RFE type
				\item \textcolor{blue}{Prompts for pre-filing documentation strategy}
			\end{itemize}
			\textbf{Checklists:}
			\begin{itemize}
				\item RFE response checklist
				\item Expert letter checklist
				\item Submission verification checklist
				\item \textcolor{blue}{Pre-filing evidence documentation checklist}
			\end{itemize}
		\end{block}
	\end{frame}
	
	\begin{frame}{Thank You - Master Your RFE Response}
		\centering
		\vspace{1cm}
		\textcolor{blue}{\Large RFE Response Toolkit}
		\vspace{0.5cm}
		\begin{itemize}
			\item Understand the 3 Prongs of Dhanasar
			\item Use LaTeX for professional responses
			\item Leverage GenAI for drafting
			\item Expert letters that address specific RFEs
			\item Respond on time, respond completely
			\item \textcolor{blue}{Document pre-filing work with timestamps}
			\item \textcolor{blue}{Pending publications confirm ongoing work, not new achievements}
		\end{itemize}
		\vspace{0.5cm}
		\textcolor{blue}{\Large Your RFE is not a denial - it's an opportunity.}
		\vspace{0.5cm}
		\textcolor{red}{This presentation is for educational purposes only. Consult a qualified immigration attorney for legal advice specific to your case.}
	\end{frame}
	
\end{document}